% STATUS: draft from Project/docs/abstract.md v0.2 (no results yet). TODO: final version with numbers (week 5).

Retrieval-Augmented Generation (RAG) is the standard way to give language models access to private documents, and running the whole pipeline on the user's own device is now practical on consumer hardware, which helps with privacy, offline use and cost. Retrieval does not, however, eliminate hallucination: generated answers still contain claims that the retrieved context does not support. Detecting such unfaithful claims as the answer is produced is a prerequisite for trustworthy local assistants. Since the check must run on every answer, it has to fit on the same constrained machine as the generator.

Automatic faithfulness verification currently follows two main paradigms. \emph{Encoder-based checkers} (NLI and fact-verification models such as AlignScore, MiniCheck and HHEM) are small and fast, but they are fixed-function classifiers: their context window is limited and they cannot take instructions or policies expressed in natural language. \emph{LLM-as-a-judge} approaches (e.g. RAGAS, ARES, Lynx) are far more flexible, but they are autoregressive. The verdict is generated token by token, so latency and cost grow with the length of the output. Small models often fail to respect the output format, and the confidence they express is rarely well calibrated. A third family has recently emerged: \textbf{System One models} (the proprietary Jev and the open re-implementations SemIf and Rizzo-Flow). These answer closed-form questions in a single forward pass, reading the probability of each answer option directly from the logits instead of generating text. The evidence for their advantage is mostly vendor-reported or small-scale, and the open implementations are general-purpose.

Three questions therefore remain open. The first is whether a \textbf{task-specialised, fine-tuned compact System One} can match the accuracy of larger generative judges and of dedicated verification encoders on a real task. The second is how much of its advantage comes from specialisation, how much from the single-pass readout, and how much from evaluating several claims over one shared context. The third is whether it delivers probabilities that are usable because they are calibrated. All three have to be answered \textbf{under the memory and latency constraints of consumer hardware} (16 GB of unified memory). No independent, controlled study appears to address them. In particular, the fairest control, a same-size LLM that emits a single answer token, is rarely included in comparisons.

This thesis designs, trains and evaluates \textbf{GroundOne}, a System One specialised in \emph{one} task: deciding whether each sentence of an answer is supported by the retrieved context. Compact open models (Qwen2.5 0.5B and 3B) are fine-tuned with LoRA on public, human-labelled faithfulness data, with a training loss restricted to the logits of the answer options. At inference, the context and the answer are encoded \textbf{once} into a KV cache. All the questions about all the answer sentences (is it supported? is it contradicted? is it acceptable under a given verification rule?) are then scored \textbf{in parallel} in a single batched pass. The probability of each answer is read directly from the logits, so there is no text generation and no output parsing, and a post-hoc temperature calibration is applied on top. Besides the plain supported/unsupported decision, GroundOne is trained to answer a three-way question and to follow verification rules written in the prompt, something a fixed-function encoder cannot do. The model is implemented on MLX and runs entirely on a 16 GB Apple Silicon laptop, where it is integrated in an end-to-end local RAG demo that flags unsupported sentences in real time.

The system is evaluated on public benchmarks with human labels (LLM-AggreFact and RAGTruth), with strict separation between training and test data and automated leakage checks. The three-way question and the verification rules are evaluated on ANLI and RAGTruth, whose labels distinguish contradicted from unmentioned content. To isolate each source of gain, GroundOne is compared against a ladder of baselines: (1) generative LLM judges of the same size, in a terse and a verbose variant, together with published results of larger models; (2) the same models with single-token readout but without fine-tuning; (3) a general-purpose System One (Rizzo-Flow); and (4) specialised verification encoders (MiniCheck and HHEM). Latency is measured as a function of the number of decisions per context. The metrics cover discrimination (balanced accuracy, AUROC), calibration (ECE, Brier score, selective-prediction curves), latency (p50/p90/p99 per answer, throughput), peak memory and output-format failure rate. All claims are stated as falsifiable hypotheses, and negative results will be reported.
